\section{Appendix: Scope and Qualifications}
This study concerns the resumed Fifth Sunday after the Epiphany in the
Vatican typical edition of the 1962 Roman Missal, under the General Roman
Calendar, on 8 November 2026. It includes eleven proper elements, without
additional commemorative orations. No local, diocesan, religious or church
calendar has been supplied. The ordinary winter Fifth Sunday and other
Masses are outside its scope. No present permission to celebrate using
the book or local schedule is determined.

The appointed Latin is the research's collated transcription, with stress
accents suppressed, ligatures resolved and line-break hyphens joined.
The printed doxology cue and abbreviated prayer conclusions are retained.
The Preface includes the dialogue and Sanctus. The English Scripture
reproduces the identified Douay--Rheims/Challoner biblical basis, whose
wording and extent differ from the adapted chants. Cummiskey's historical
English supplies the three orations and Preface, including his own
abbreviations and same-witness expansions. Neither English witness is an
approved proclaimed translation of the 1962 Missal.

The two interpretations are editorial syntheses of this Mass. The Fathers
and saints are credited with their cited passage explanations and doctrinal
arguments. The literal exiles, the Gospel's world and the Christian
assembly are distinguished. Chrysostom's immediate identification of tares
with heretics and Augustine's broader address to wicked Christians remain
attributed. Chrysostom's permission of nonlethal suppression qualifies his
prohibition of killing. These passages do not settle every question of
civil force or ecclesiastical discipline. Fromage's continuation corroborates
ten shared texts without developed exposition or the Trinity Preface.\footnote{Fromage's continuation of \work{The Liturgical Year}, trans. Shepherd, XI (Duffy, 1883), appendix, Fifth Sunday after the Epiphany, pp.~524--528.}

The research used historical English patristic translations, retained
Latin transcriptions and the specified printed-page witnesses. No new
Greek collation is claimed. Bellarmine's English is O'Sullivan's abridgment;
Bede's Mark exposition is used from its retained transcription, with no
new printed-page collation. Aquinas's checked Colossians page covers
3:16--17; the preceding missing scan pages supply no claim. The bounded
reception sweep, uninspected leads, exact evidence states and source limits
remain in \texttt{research/scope.md}, \texttt{research/interpretations.md},
\texttt{research/chronology-notes.md}, \texttt{research/source-bindings.toml}
and \texttt{propers/verified.md}. The historical appendix retains the
distinction between traditional attributions and modern critical proposals.

The inherited Latin and historical English retain their recorded United
States public-domain bases through the identified antecedents. The modern
1962 facsimile is a collation source; its whole scan and modern material
are not cleared for redistribution. Modern biblical introductions are
consultation sources, without reproduced introduction bodies. The original
commentary is distinct from received prayer. Author proof verification is
recorded separately from the independent reviews of the three documents;
no external or ecclesiastical approval is claimed.

\section{References}
\subsection{Liturgical and biblical texts}
\begin{itemize}
\item \work{Missale Romanum}, Vatican editio typica (1962),
\href{https://media.churchmusicassociation.org/pdf/missale62.pdf}{CMAA
facsimile}, pp.~420--421, nos.~1749--1758; Trinity Preface,
p.~293, no.~1082. General Roman Calendar; the researched appointment
and occurrence are recorded in \texttt{research/context.md}.
\item \work{Missale Romanum}, Pustet, Ratisbon (1862),
pp.~42--44, 357--358, Latin antecedents; Trinity opening and clause,
pp.~239--240. Eugene Cummiskey, \work{The Roman Missal translated into
the English language for the use of the laity}, first revised edition
(Philadelphia, 1861), own Fifth-Sunday prayers, pp.~101--102; Preface dialogue, opening,
Trinity clause and Sanctus, pp.~xxvii--xxix. Historical wording retains
its independent public-domain status.
\item \work{Douay--Rheims Bible}, Challoner revision, composite electronic
witness \href{https://www.gutenberg.org/ebooks/1581}{Gutenberg 1581}:
Jer 29, 36; Ps 43, 84, 129 (Vulgate numbering); Col 1:2; 3:1--17; 4:7--18;
Matt 13; Mark 11. The appointed selections and complete literary context
are distinguished in the commentary.
\end{itemize}
\subsection{Exegesis and theology}
\begin{itemize}
\item John Chrysostom, \work{Homilies on Matthew}, XLVI.1,
NPNF first series 10, retained CCEL historical English, revised by
George Prevost and M.~B. Riddle. \work{Homilies on Colossians},
VIII on 3:12--15 and IX opening on 3:16--17, NPNF first series 13,
retained CCEL historical English. Translation lemmata remain distinct
from the Missal Latin.
\item Augustine, Sermon XXIII [Benedictine LXXIII], §§1--4,
NPNF first series 6, retained CCEL historical English.
\work{Expositions of the Psalms}, Ps 43 §§1,3--8;
Ps 84 §§2--6; Ps 129 §§1--8, NPNF first series 8,
retained CCEL text under headings Psalms 44, 85 and 130 respectively.
These are the historical English witnesses, without a fresh original-Latin
collation.
\item Jerome, \work{Commentary on Jeremias}, V, on 29:10--14,
PL 24 (Paris, 1865), cols.~893--894; surrounding context
cols.~895--896. Original-page witness checked in research.
\item Robert Bellarmine, \work{Commentary on the Book of Psalms},
Ps 129, verses 1--2, John O'Sullivan's abridged historical English
(1866), retained text from eCatholic2000.
\item Bede, \work{In Marci evangelium expositio}, on Mark 11:22--26,
paragraphs \textit{Et respondens Jesus} and \textit{Et cum stabitis ad
orandum}, retained Wikisource/Migne Latin transcription. Its present
fragment-admission limit is recorded in the research.
\item Thomas Aquinas, \work{Summa theologiae} II--II, q.~83,
aa.~15--16, English Dominican translation, Gutenberg 18755.
\work{Commentary on Colossians}, chapter 3, lectio 3,
on 3:16--17, \work{In omnes Pauli epistolas commentaria},
II (Dessain, 1857), p.~442.
\end{itemize}
The historical appendix gives the exact articles, official introductions
and AAS responses used for its chronological and geographical claims.
All online witnesses were inspected in the research completed on
2 November 2026; their exact identities and source-local limits remain
in the binding record.
